% Propietario conservado: /Users/ruben/Documents/New project/output/EXTREMA_MEDIA_RAZON_RECIPROCIDAD_ES_EN_20260918/ARTICULO/ES/source/nuclear/04.tex
\chapter{Incidence exceptionnelle et transport des formes}
\label{x_reciprocidad:nuclear:4}
\section{Génération du registre et conservation de l’incidence}

\subsection{Registre dodécaphasique de l’histoire enrichie}

La section~\ref{x_reciprocidad:subsec:k-registro-integral} définit

\[
\mathcal R_{12}(x)=((E_m,\tau_m,\sigma_m,\kappa_m,\Lambda_m))_{m=1}^{12}.
\]

Il conserve la fenêtre, la sous-route, l’orientation, la retenue et le livre des incidences. La section~\ref{x_reciprocidad:subsec:k-terminal} fixe

\[
x_{\rm term}=(\mathcal U_{108}\cdots\mathcal U_1)x_{\rm seed},
\quad \mathcal U_t=\mathsf{Upd}_t\mathsf{Tra}_t\mathsf{Sel}_t,
\]

\[
R_{\rm sgn}=\Pi_H\mathcal E_{108}^{90,120}\mathcal R_{12}:
\mathcal X_{\rm TPK}^{\rm term,mem}\to\mathbb Z^{12}.
\]

La reconstruction algébrique suivante commence à cette sortie typée. Elle ne remplace pas l’évaluation des événements antérieurs et n’obtient pas leurs coordonnées en imposant \(K\) ou alpha comme cible.

\subsection{Normalisation entière et normalisations isométriques}

Pour les orbites de pas trois \((1,4,7,10)\), \((2,5,8,11)\), \((3,6,9,12)\), le lemme~\ref{x_reciprocidad:lem:k-hadamard} fixe la matrice symétrique \(H=H_{12}\) et prouve \(H^2=4I\).

\[
\mathcal L_H=\{u\in\mathbb Z^{12}:Hu\equiv0\pmod4\},
\qquad \mathscr K(u)=\tfrac14Hu,
\qquad \mathscr K^{-1}(k)=Hk.
\]

C’est un isomorphisme entier entre le sous-réseau des registres signés et \(\mathbb Z^{12}\). Le facteur \(1/4\) réalise l’inversion, et non une isométrie euclidienne.

Comme conséquence immédiate de la même identité, \(J_H=H/2\) satisfait \(J_H^*J_H=I\) : c’est le transport isométrique réel. Si \(u=Hk\), on a \(u/2=J_Hk\) et \(\|u\|^2=4\|k\|^2\). Confondre \(H/4\) avec \(H/2\) changerait le bilan d’un facteur quatre.

Dans \cref{x_reciprocidad:exc:entrelazador,x_reciprocidad:exc:isometria}, le plan combinatoire construit détermine

\[
I:\mathbb R^{12}\to\mathbb R^{132},\qquad
(Iv)(H')=\sum_{p\in H'}v_p,\qquad I^*I=36I_{12}+30J_{12}.
\]

Les coefficients proviennent des incidences : chaque point appartient à 66 hexades et chaque paire à 30. Sur \(V=\mathbf1^\perp\), \(J_{12}\) disparaît, de sorte que \(U_W=I/6:V\to E_{\rm inc}\) est une isométrie. Le facteur \(1/6\) est le normalisateur de cet écran, non l’inverseur de \(K\).

\subsection{Raffinement des préfixes et partitions d’événements}

La section~\ref{x_reciprocidad:exc:doble-lectura} construit des repères causaux \(f_{j+1}=g_j(x)f_j\) et des codages \(Q_n\) du préfixe. Elle prouve

\[
r_{n,m}Q_n=Q_mp_{n,m},\qquad
\gamma_{A^f}(wf^{-1})=\gamma_A(w)(f^{-1}\oplus f^{-1}).
\]

La preuve utilise le fait qu’un prolongement conserve les mots, les transports et les repères précédents. La limite projective produit \(Q_\infty\). La récupération du mot visible est due à la première composante de \(\gamma\) ; le texte distingue expressément cette portée de la récupération de toute mémoire profonde et du passage aux supports.

La composition additive et affine du registre, décrite dans \cref{x_reciprocidad:subsec:k-registro-integral,x_reciprocidad:iv:cont:concatenacion}, prouve une seconde conservation : avec des contributions prospectives compatibles \((\delta w_t,\delta q_t)\), l’inversion du registre est additive. Concaténer des segments ou subdiviser la même liste d’événements donne le même registre final. La version avec transport non trivial utilise le cocycle affine. Conserver le total lors d’un changement de partition n’équivaut pas à affirmer que le registre cesse de croître lorsqu’on ajoute de nouveaux événements.

\section{Transport K–Witt des normes, des opérateurs et des bilans}\label{x_reciprocidad:nuc04:kwitt}

On réunit ici une formalisation des résultats récupérés, avec la preuve algébrique qui suit. Cet assemblage n’attribue aucune priorité nouvelle à ses identités de départ.

Soient \(P_0=I-J_{12}/12\), \(V=P_0\mathbb R^{12}\), \(W_H=J_HV\) et

\[
B:=U_WJ_H|_{W_H}:W_H\longrightarrow E_{\rm inc}.
\]

Comme \(J_H^2=I\), c’est un isomorphisme isométrique entre espaces de dimension onze. Pour un registre \(k\), sa coordonnée signée normalisée est \(J_Hk\). Si \(\widehat P_0=J_HP_0J_H\),

\[
B\widehat P_0(J_Hk)=U_WP_0k,
\]

\[
\|\widehat P_0(J_Hk)\|^2
=\|U_WP_0k\|^2
=\sum_{p=1}^{12}k_p^2-\frac{(\sum_pk_p)^2}{12}.
\]

Aucune valeur de constante n’a été introduite pour obtenir cette égalité : elle provient de \(H^2\) et du Gramien d’incidence. La moyenne éliminée doit être conservée séparément si l’on veut retrouver tout \(k\).

Pour \(g\in\operatorname{Aut}(\mathcal H)\), soit \(\rho_{12}(g)\) sa permutation des points et \(\rho_{\rm hex}(g)\) celle des hexades. L’action sur les coordonnées signées est

\[
\widehat\rho_H(g)=J_H\rho_{12}(g)J_H
=\tfrac14H\rho_{12}(g)H.
\]

Il est orthogonal, conserve \(W_H\) et conserve également le sous-réseau \(\mathcal L_H\) : si \(u=Hk\), son image est \(H\rho_{12}(g)k\). L’entrelacement matériel de \cref{x_reciprocidad:exc:entrelazador} donne

\[
B\widehat\rho_H(g)=\rho_{\rm hex}(g)B.
\]

\begin{proof}
substituer les définitions, simplifier \(J_H^2\) et appliquer \(U_W\rho_{12}(g)=\rho_{\rm hex}(g)U_W\). Le registre entier, le produit scalaire centré et l’action incidencielle sont ainsi préservés, chacun dans son domaine.
\end{proof}

Pour un opérateur \(A\) déjà construit sur \(W_H\), on obtient \(A_{\rm inc}=BAB^*\). Par substitution,

\[
BA=A_{\rm inc}B,\quad
\langle z,Az\rangle=\langle Bz,A_{\rm inc}Bz\rangle.
\]

Par conséquent, si un transport \(C\) sur \(W_H\) conserve \(A\), \(C^*AC=A\), son transport \(C_{\rm inc}=BCB^*\) conserve \(A_{\rm inc}\). La condition selon laquelle un \(C\) particulier est une publication de la dynamique TPK reste dans la composition causale qui le produit. L’isométrie ne transforme pas à elle seule une permutation d’incidence en cycle nonadique.


Il n’est pas nécessaire d’imposer la commutation pour comparer des cartes transportées. Si \(R\) est unitaire, \(C'=RCR^*\) et \(T(C)=(8I+C)/9\), alors \(T(C')R=RT(C)\) par substitution. Avec \(A'=RAR^*\), le bilan dans une carte se conjugue à celui de l’autre. La commutation \([R,C]=0\) ne correspond qu’au problème supplémentaire consistant à maintenir fixe cette même couture \(C\). Cette distinction coïncide avec la naturalité réticulaire de la section~\ref{x_reciprocidad:nuc04:c3}.

Une action finie concrète disponible est l’étoile sélectionnée par le drapeau :

\[
\sigma_{B_K}=(1\ 3)(2\ 11)(5\ 7)(8\ 12).
\]

La section~\ref{x_reciprocidad:exc:estrella} prouve sa construction et la distingue du groupe engendré par les 495 étoiles. L’étoile individuelle donne \(C_2\) ; la famille donne \(M_{12}\). Transporter la marque et l’origine donne la covariance entre repères. Si l’on vise une symétrie d’un repère fixe, le domaine approprié est son stabilisateur, et non automatiquement \(M_{12}\) entier. Si elle doit en outre commuter avec \(C\), on utilise le sous-groupe qui satisfait l’équation concrète \([\widehat\rho_H(g),C]=0\).

\section{Action d’ordre trois sur les fibres réticulaires}\label{x_reciprocidad:nuc04:c3}

La construction de \cref{x_reciprocidad:nuclear:fibras-app-witt,x_reciprocidad:nuclear:alineamiento-witt} produit douze fibres \(A_2\) étiquetées par \((i,\mu,\varepsilon)\in\mathbb Z/3\mathbb Z\times\{0,1\}\times\{0,1\}\). L’alignement \(\lambda(i,\mu,\varepsilon)=4i+2\mu+\varepsilon\pmod{12}\) est bijectif. Si \(C_{A_2}\) est le Coxeter et \(R\) la réflexion telle que \(RC_{A_2}R=C_{A_2}^{-1}\), les trivialisations et les repères transportés fixent

\[
g_b=\iota_b^{-1}P_bC_{A_2}P_b^{-1}\iota_b,
\quad T_{i\mu\varepsilon}=\iota_{\lambda(i,\mu,\varepsilon)}^{-1}P_{\lambda(i,\mu,\varepsilon)}R^\mu,
\]

\[
g_\lambda T_{i\mu\varepsilon}
=T_{i\mu\varepsilon}C_{A_2}^{(-1)^\mu}.
\]

La somme directe \(T_\lambda:W_{\rm APP}\otimes\mathbb C\to\bigoplus_{b=1}^{12}(A_{2,b}\otimes\mathbb C)\) est un isomorphisme \(C_3\)-équivariant. La preuve consiste en la conjugaison et en la bijection des douze étiquettes : elle conserve la paire, le feuillet et l’orbite. Ce n’est pas une coïncidence de dimensions.

Les constructions de \cref{x_reciprocidad:km:palabra-total,x_reciprocidad:km:estabilidad-vecino} utilisent \(q^\star=(1,1,1,1,1,1)\) et son mot \(c^\star=(q^\star,q^\star A_W)\) à support total pour orienter les douze blocs. Le \(g_c\) résultant est d’ordre trois, n’a pas de vecteurs fixes non nuls et préserve le recollement de Niemeier. Pour le vecteur incidenciel \(v\) de norme quadratique \(54\),

\[
N_0=\ker(\langle\cdot,v\rangle\bmod3),\qquad
N_v=N_0+\mathbb Z\frac v3,
\]

\[
y_*=(g_cv-v)/3\in N_0,\qquad
z_*=(g_c^{-1}v-v)/3\in N_0.
\]

La preuve identifie leurs classes discriminantes comme \(c^\star\) et \(-c^\star\), et leurs produits avec \(v\) comme \(-27\). D’où \(g_cN_0=N_0\) et \(g_c(v/3)=v/3+y_*\), puis \(g_cN_v=N_v\). La structure conservée comprend le recollement, le noyau de congruence et la classe qui engendre le voisin de Leech.

Le prolongement aux 24 orientations de support total et l’élévation réticulaire d’un automorphisme monomial \(h\) du code sont développés dans la section~\ref{x_reciprocidad:paper:24-orientaciones}. Avec toutes les marques transportées,

\[
g_{hu}=\widetilde h\,g_u\,\widetilde h^{-1}.
\]

Ce \(C_3\) est un groupe permis réellement entrelacé avec les fibres APP et la structure exceptionnelle. On ne l’identifie ici ni à \(C_9\) ni à son quotient par une simple relation d’ordres. L’élévation au réseau n’est pas non plus remplacée par la matrice ternaire sans son action sur les fibres \(A_2\).

\section{Raffinement de l’énergie et de sa différentielle}
\label{x_reciprocidad:paper:refinamiento-energia}

L’opérateur d’énergie de \eqref{x_reciprocidad:vii:eq:energia-memoria-discreta} est de la forme \(D_{9,m}^*W_mD_{9,m}\), avec \(U_a\) transportant l’orientation et la retenue. Les domaines de \cref{x_reciprocidad:vii:subsec:funcional-memoria} sont :

\[
\mathcal H_m=\bigoplus_vH_v,\quad
\mathcal K_m=\bigoplus_{a\in\mathcal E_m^+}H_{t(a)},\quad
(D_mf)_a=f_{t(a)}-U_af_{s(a)},\quad
E_m=\langle D_mf,\mathsf W_mD_mf\rangle.
\]

Les applications fixes \(J_m\) et \(K_m\) satisfont, pour la famille \(C^1\) considérée,


\[
D_{m+1}(t)J_m=K_mD_m(t),\qquad
K_m^*\mathsf W_{m+1}(t)K_m=\mathsf W_m(t).
\]

Le \cref{x_reciprocidad:vii:thm:refinamiento-respuesta} prouve \(E_{m+1}(J_mf)=E_m(f)\) par substitution, et l’égalité des différentielles s’obtient en différentiant les deux identités. Pour un raffinement dépendant du paramètre, on conserve également le terme \(D_{m+1}\dot J_mf\).

Avec \(v_a=U_af_s\) et \(q=\mathsf W_mD_mf\), la réponse à \(\dot U_a=A_aU_a\) est

\[
\mathfrak j_m(A)=-2\operatorname{Re}\sum_a\langle q_a,A_av_a\rangle,
\quad \mathcal J_a=v_aq_a^*-q_av_a^*.
\]

Le produit de transports \(U_\gamma=U_N\cdots U_1\) vérifie

\[
\dot U_\gamma U_\gamma^{-1}
=\sum_{k=1}^N B_kA_kB_k^{-1},\qquad B_k=U_N\cdots U_{k+1}.
\]

Cette formule retient l’ordre de chaque incidence et transporte son covecteur de manière contragrédiente. Le corollaire~\ref{x_reciprocidad:vii:cor:refinamiento-corriente-conexion} compose cette différentielle avec les variations de connexion et prouve, pour les raffinements compatibles de l’action,

\[
s_m=P_m^{\mathsf T}s_{m+1},\qquad
M_m\sigma_m=P_m^{\mathsf T}M_{m+1}\sigma_{m+1}.
\]

Ici \(P_m\) porte des variations de connexion et \(M_m\) est l’appariement cellulaire ; ce ne sont ni le projecteur de moyenne nulle ni la matrice de Hadamard.

\subsection{Assemblage avec l’incidence}

Dans une réalisation de fibres \(W_H\) dont les actions de repère sont \(\widehat\rho_H(g_a)\), la conjugaison par \(B\) envoie chaque transport, champ et poids sur \(E_{\rm inc}\). D’après l’entrelacement démontré dans la section~\ref{x_reciprocidad:nuc04:kwitt}, les opérateurs de différences satisfont \(D_{\rm inc}B_{\rm vertices}=B_{\rm edges}D_H\). En transportant également \(J_m\), \(K_m\) et \(\mathsf W_m\), les deux carrés de raffinement précédents restent les mêmes carrés conjugués. Par substitution, ils conservent l’énergie ; pour \(B\) fixe, ils conservent en outre la différentielle et les covecteurs.

La loi de bilan et ses covecteurs sont transportés vers la représentation exceptionnelle en conservant leur compatibilité avec le raffinement. Chaque application utilise les transports de la dynamique spécifiée et les sous-espaces qu’ils préservent. Le domaine des variations est défini dans \cref{x_reciprocidad:vii:subsec:funcional-memoria}.
